\definemathmatrix[pmatrix]        % defining matrix with parentheses
	[matrix:parentheses]
	[simplecommand=pmatrix]

\margintext{2024-10-30}
\section{Bernoulli Formel}

\startformula
(\startpmatrix
	   \NC n \NR
	   \NC k \NR
\stoppmatrix) = \frac{n!}{k! \cdot (n-k)!}
\stopformula

\blank
\hrule
\blank

\startitemize[packed]
\item 20 aufgaben multiple choice
\item jede aufgabe hat genau 5 aufgaben, genau eine ist richtig
\item Jede aufgabe wird mit 1 BE berechnet
\stopitemize

Wie wahrscheinlich,  alle richtig: $(\frac{1}{5})^2$

Wie wahrscheinlich genau 50\% richtig

\startformula
(\startpmatrix
	   \NC 20 \NR
	   \NC 10 \NR
\stoppmatrix) \cdot 0.2^{10} \cdot 0.8^{10}
\stopformula

{\bfa Eine Note 1 (mindestens 80\% richtig, 17 Aufgaben)}

\startformula
(\startpmatrix
	   \NC 20 \NR
	   \NC 17 \NR
\stoppmatrix) \cdot 0.2^{17} \cdot 0.8^{3}
\stopformula

\startformula
\sum^{20}_{n=17} 
(\startpmatrix
	   \NC 20 \NR
	   \NC n \NR
\stoppmatrix) \cdot 0.2^{n} \cdot 0.8^{20 - n} = \frac{76081}{95367431640625} \approx 7.978 \cdot 10^{-10}
\stopformula

\blank
{\bfa Einen unterkurs (<45\%)}
\blank

\margintext{
{\bf TR-Befehle:}\crlf
Der C.A.S. befehl ist: {\em binomialPDf} fuer einfache Wahrscheinlichkeiten und {\em binomialCDf} fuer kommulierte Wahrscheinlichkeiten

Diese gehören {\bf nicht} die Dokumentation. In diese gehört zum Beispiel diese hier:

\startformula
P(9 \leq x \leq 12) \approx 9,97 \cdot 10^{-3}
\stopformula
}

\startformula
P(x\leq 8)=
\sum^{8}_{n=0} = P(x=n)
\sum^{8}_{n=0} 
(\startpmatrix
	   \NC 20 \NR
	   \NC n \NR
\stoppmatrix) \cdot 0.2^{n} \cdot 0.8^{20 - n} = \frac{94415494316032}{95367431640625} \approx 0.99
\stopformula

\startdefinition[title={Kommulieren}]
Aufaddieren von Wahrscheinlichkeiten heist "kommulieren"
\stopdefinition
